\section{Schubert classes and the attachment modulo six}\label{sec:schubert}

Work first in the positive chamber, with the actual extension
\eqref{eq:actual-extension}. Write
\(h=c_1(\mathcal U)\), \(p=c_2(\mathcal U)\).
The integral Schubert presentation \cite[Section~3]{buch} is
\begin{equation}\label{eq:schubert-ring}
 H^*(G,\Z)=\Z[h,p]/(h^3-2hp,\ h^2p-p^2),\qquad
 |h|=2,\quad |p|=4,
\end{equation}
with
\begin{equation}\label{eq:schubert-evaluations}
 \int_G h^4=2,\qquad\int_G h^2p=\int_G p^2=1.
\end{equation}
The second degree-four Schubert class is \(q=h^2-p\);
thus \(\int_G q^2=1\) and \(\int_G pq=0\).
These are the classical two Schubert directions of \(\Gr(2,4)\).
They are represented by planes of the two types: planes containing
a fixed line, and planes contained in a fixed three-space.

\subsection{Removing the universal-family twist}
In the lattice embedding \eqref{eq:rankseven-lattice}, put
\(D=e-22f\). The integral Mukai vectors
\[
 x=(0,-f,0),\qquad y=(0,3f,1)
\]
are topological \(K\)-classes; algebraicity is not required.
Their Euler pairings are
\[
 \chi(A,x)=1,\quad\chi(B,x)=0,\qquad
 \chi(A,y)=0,\quad\chi(B,y)=1.
\]
If the restricted universal family differs from
\eqref{eq:actual-extension} by a line \(L\), then
\(I_x|_G=\mathcal U\otimes L\) and \(I_y|_G=L\).
The global virtual rank of \(I_y\) is one. Tensor the universal
class by \(\det(I_y)^{-1}\); this removes \(L\) on \(G\).
For the resulting normalized indices,
\begin{equation}\label{eq:index-restrictions}
 I_z|_G=m_z[\mathcal U]+n_z[1],\qquad
 c(I_z|_G)=(1+h+p)^{m_z},\qquad
 m_z=\chi(A,z),\quad n_z=\chi(B,z).
\end{equation}
The binomial expansion is integral for every integer \(m_z\),
including negative ones.

In particular \(c_1(I_x)\) and \(c_2(I_x)\) extend \(h,p\)
integrally to \(M\). Since \(h,p\) generate
\eqref{eq:schubert-ring}, restriction is onto in every degree.
Duality makes every corresponding inclusion
\(H_k(G,\Z)\to H_k(M,\Z)\) primitive.
The class \([G]\) is primitive because the integral extension
\(c_2(I_x)^2\) evaluates to one on it.
This is stronger than the information from divisors:
every divisor restricts to a multiple of \(h\), whose fourth
power has evaluation two.

\subsection{The boundary and complement}
The Plücker embedding identifies \(G\) with a smooth quadric
fourfold in \(\PP^5\). The Euler and normal sequences give
\[
 c(TG)=\frac{(1+h)^6}{1+2h},\qquad
 c_4(TG)=3h^4=6p^2 .
\]
Since \(G\) is Lagrangian, its normal bundle is \(T^*G\), so
\begin{equation}\label{eq:self-intersection}
 [G]^2=\int_G c_4(T^*G)=6 .
\end{equation}
For the boundary \(L=S(T^*G)\), the sphere-bundle Gysin sequence
and \eqref{eq:schubert-ring} give
\begin{equation}\label{eq:boundary-groups}
 H^j(L,\Z)=
 \begin{cases}
 \Z,&j=0,2,6,9,13,15,\\
 \Z^2,&j=4,11,\\
 \Z/6,&j=8,\\
 0,&\text{otherwise}.
 \end{cases}
\end{equation}
The only nonzero Euler multiplication is from degree zero
to degree eight, where it is multiplication by six.
This is a smooth boundary of the resolved neighborhood;
the target has a six-dimensional singular stratum through \(o\),
so it is not the link of an isolated singularity.

Put \(V=M\setminus G\). Markman's theorem gives torsion-free even
cohomology of \(M\), and the primitive inclusions just proved
make the relative sequence torsion-free. More explicitly,
\(H^k(M,V;\Z)=H_{16-k}(G,\Z)\); its map to \(H^k(M,\Z)\)
is the dual of restriction and is a primitive injection whenever
nonzero. Odd groups of \(V\) vanish. The classical \(X^{[4]}\)
Betti numbers, obtained from Göttsche's formula \cite{gottsche},
therefore give
\[
\begin{array}{c|rrrrrrrrr}
k&0&2&4&6&8&10&12&14&16\\ \hline
\rank H^k(M)&1&23&300&2852&19298&2852&300&23&1\\
\rank H^k(V)&1&23&300&2852&19297&2851&298&22&0 .
\end{array}
\]
The ranks removed in degrees \(8,10,12,14,16\) are
\(1,1,2,1,1\). These ranks do not by themselves give the
marked complement ring needed for attachment.

\subsection{The calculated congruence}
Apply Theorem~\ref{thm:attachment} with \(n=8\),
\(\Lambda=\Z[G]\) and \(B=(6)\).
For a marked universal-index Chern monomial of degree eight
on \(V\), substitute \eqref{eq:index-restrictions}, reduce by
\eqref{eq:schubert-ring}, and evaluate using
\eqref{eq:schubert-evaluations}, modulo six.
This specifies a homomorphism
\(\kappa:H^8(V,\Z)\to\Z/6\) on all integral generators.
In detail, for \(m=m_z\),
\begin{equation}\label{eq:chern-polynomials}
 \begin{aligned}
 c_1(I_z)|_G&=mh,\\
 c_2(I_z)|_G&=mp+\binom m2h^2,\\
 c_3(I_z)|_G&=\frac{m(m^2-1)}3\,hp,\\
 c_4(I_z)|_G&=\frac{m^2(m^2-1)}{12}\,p^2 .
 \end{aligned}
\end{equation}
The displayed coefficients are integer-valued for every integer
\(m\). The ring relations give the rule for all products, including
products involving different indices. The all-relations argument
in Theorem~\ref{thm:attachment} proves that this calculation
depends only on the specified complement class.

We obtain the explicit integral reconstruction
\begin{equation}\label{eq:mod-six}
 H^8(M,\Z)\simeq
 \{(u,N)\in H^8(V,\Z)\oplus\Z:
                  N\equiv\kappa(u)\pmod6\}.
\end{equation}
Moreover
\(\kappa(c_2(I_x)^2|_V)=1\), so \(\kappa\) is onto.
The higher index class supplies the odd evaluation which fourth
powers of divisor classes cannot supply. Both Schubert directions,
and their integral attachment, are thus visible in the same
compact moduli family.

For the negative chamber the pure extension axis instead restricts
the normalized universal class to \(A\boxtimes\mathcal S+B\),
where \(\mathcal S\) is the tautological subbundle in the
image-of-\(Y_1\) description. Its Chern classes
\(c_1(\mathcal S)=-h\), \(c_2(\mathcal S)=q=h^2-p\)
again generate the full integral ring. The same duality argument
proves all primitive inclusions in that chamber as well.

The core is also genuinely different from a toric core component.
Indeed \(\Pic(G)=\Z h\), so a connected automorphism preserves the
Plücker embedding and acts through the projective orthogonal group
of its smooth quadric. Its maximal torus has dimension three.
A dense torus action on a toric fourfold would require dimension
four. This observation is not needed for \eqref{eq:mod-six};
the two independent Schubert classes already explain the extra
integral information.
